\abstract{РЕФЕРАТ}

Объем работы равен \formbytotal{lastpage}{страниц}{е}{ам}{ам}. Работа содержит \formbytotal{figurecnt}{иллюстраци}{ю}{и}{й}, \formbytotal{tablecnt}{таблиц}{у}{ы}{}, \arabic{bibcount} библиографических источников и \formbytotal{числоПлакатов}{лист}{}{а}{ов} графического материала. Количество приложений – 2. Графический материал представлен в приложении А. Фрагменты исходного кода представлены в приложении Б.

Перечень ключевых слов: интерактивное приложение, компонентный подход, компонентно-ориентированная система, Java, LWJGL, OpenGL, GLSL, сцена, сущность, компонент, система обработки, рендеринг, физическая подсистема.

Объектом разработки является компонентно-ориентированная программная система для построения интерактивных приложений на языке Java.

Целью выпускной квалификационной работы является разработка программной системы, предоставляющей Java-разработчику основные средства создания интерактивных приложений: управление приложением и сценами, компонентную модель объектов, обработку пользовательского ввода, графический вывод, выполнение пользовательских скриптов и базовую физическую обработку.

В процессе выполнения работы проведен анализ предметной области и существующих программных средств построения интерактивных приложений. Сформулированы требования к программной системе, спроектирована архитектура приложения, сцены, сущностей, компонентов и систем обработки. В качестве основных технологий использованы язык Java, библиотека LWJGL, графический API OpenGL, язык шейдеров GLSL, а также библиотеки GLFW и STB.

В результате работы реализованы ядро компонентной модели, управление жизненным циклом приложения и сцен, графическая подсистема, подсистема ресурсов, обработка пользовательского ввода, система пользовательских скриптов, отладочная визуализация и базовая физическая подсистема. Для проверки работоспособности разработаны демонстрационные сцены, показывающие использование камеры, освещения, текстур, материалов, скриптов, пользовательского ввода, коллайдеров и физических тел.

Разработанная программная система может использоваться как основа для создания учебных, демонстрационных и исследовательских интерактивных приложений на Java, а также для изучения внутреннего устройства компонентной архитектуры, графической подсистемы и базовой физической обработки.

\selectlanguage{english}
\abstract{ABSTRACT}
  
The volume of work is \formbytotal{lastpage}{page}{}{s}{s}. The work contains \formbytotal{figurecnt}{illustration}{}{s}{s}, \formbytotal{tablecnt}{table}{}{s}{s}, \arabic{bibcount} bibliographic sources and \formbytotal{числоПлакатов}{sheet}{}{s}{s} of graphic material. The number of applications is 2. The graphic material is presented in annex A. Source code fragments are presented in annex B.

List of keywords: interactive application, component-based approach, component-oriented system, Java, LWJGL, OpenGL, GLSL, scene, entity, component, processing system, rendering, physics subsystem.

The object of the development is a component-oriented software system for building interactive applications in Java.

The purpose of the final qualification work is to develop a software system that provides a Java developer with the basic tools for creating interactive applications: application and scene management, a component-based object model, user input processing, graphical output, user script execution and basic physics processing.

During the work, the subject area and existing software tools for building interactive applications were analyzed. The requirements for the software system were formulated, and the architecture of the application, scene, entities, components and processing systems was designed. The main technologies used in the development are the Java programming language, the LWJGL library, the OpenGL graphics API, the GLSL shading language, as well as the GLFW and STB libraries.

As a result of the work, the core of the component model, the life cycle management of the application and scenes, the graphics subsystem, the resource subsystem, user input processing, the user script system, debug visualization and the basic physics subsystem were implemented. Demonstration scenes were developed to verify the functionality of the system and show the use of camera, lighting, textures, materials, scripts, user input, colliders and physical bodies.

The developed software system can be used as a basis for creating educational, demonstration and research interactive applications in Java, as well as for studying the internal structure of component architecture, the graphics subsystem and basic physics processing.
\selectlanguage{russian}
